\chapter{Module Theory}

\section{Basic Definitions}

\begin{definition}
    Let $R$ be a Ring. A \textit{left Module over $R$} is an Abelain Group $(M, +)$ with: \\
        An action $\cdot$ of $R$ on $M$ satisfies: for any $r,s \in R,\ m,n \in M$,
        $\begin{cases}
            (r + s) \cdot m = r \cdot m + s \cdot m \\
            (r s) \cdot m = r \cdot (s \cdot m) \\
            r \cdot (m + n) = r \cdot m + r \cdot n 
        \end{cases}$. \\
    In addition, if $R$ has an identity $1$, $1 \cdot m = m$
\end{definition}

\begin{definition}
    Let $M$ be a left Module over a Ring $R$. \\
    A subgroup $N \leq M$ is called \textit{submodule} of $M$ if: which is closed under the action of $R$ on $M$. \\
    That is, for all $r \in R$, $n \in N$, $r \cdot n \in N$.
\end{definition}

\begin{definition}
    Let $M, N$ be lef Modules over a Ring $R$. \\
    A map $\varphi:M \to N$ is called \textit{Module-Homomorphism} if: for any $x,y \in M, r \in R$,
    $\begin{cases}
        \varphi(x + y) = \varphi(x) + \varphi(y) \\
        \varphi(r \cdot x) = r \cdot \varphi(x)
    \end{cases}$.
\end{definition}

\begin{lemma}
    \textbf{Submodule Criteria} \\
    Let $M$ be a left Module over a Ring $R$. A subgroup $N \leq M$ is a submodule of $M$ if and only if
    \begin{align*}
        N \neq \emptyset,\
        \forall x, y \in N,\ r \in R,\ r \cdot x + y \in N
    \end{align*}
\end{lemma}


\begin{definition}
    Let $R$ be a Ring with identity. \\
    A Ring $A$ with identity is called \textit{$R$-Algebra} if: there exists a Ring Homomorphism $f:R \to A$ such that
    \begin{align*}
        f(1_R) = 1_A,\ Z(A) \subseteq f[R]
    \end{align*}
    Remark: $Z(A) \defequal \{r \in A \mid \forall a \in A,\ ar = ra\}$.
\end{definition}

\begin{definition}
    Let $A, B$ be $R$-Algebras. \\
    A Ring-Homomorphism $\varphi:A \to B$ is called \textit{Algebra-Homomorphism} if: $\varphi(1_A) = 1_B$ and
    \begin{align*}
        \forall r \in R, a \in A,\ \varphi(r \cdot a) = r \cdot \varphi(a)
    \end{align*}
\end{definition}

\newpage
\section{Module of Module-Homomorphism}
\begin{definition}
    Let $M,N$ be left Modules over a Ring $R$. \\
    Define two operation on $\operatorname{Hom}_R(M,N)$: for any $\varphi, \psi \in \operatorname{Hom}_R(M,N), r \in R$,
    \begin{align*}
        \begin{cases}
            \varphi + \psi:M \to N:m \mapsto \varphi(m) + \psi(n) \\
            r \cdot \varphi: M \to N : m \mapsto r \cdot \varphi(m)
        \end{cases}
    \end{align*}
    Then, $(\operatorname{Hom}_R(M,N), +, \ \cdot\ , R)$ becomes a left Module over $R$.
\end{definition}

\newpage
\section{Direct Sum of Modules}
\begin{definition}
    Let $M$ be a left Module over a Ring $R$ and let $N_1, \ldots, N_n$ be submodules of $M$. \\
    Define \textit{Sum} of $N_i$:
    \begin{align*}
        N_1 + N_2 + \cdots + N_n \defequal \{a_1 + a_2 + \cdots + a_n \mid a_i \in N_i,\ \forall 1 \leq i \leq n \}
    \end{align*}
    Let $A$ be any subset of $M$. Define \textit{submodule generated by $A$}:
    \begin{align*}
        RA \defequal \{r_1 \cdot a_1 + \cdots + r_n \cdot a_n \mid n \in \mathbb{N}, r_i \in R, a_i \in A\}
    \end{align*}
    A submodule $N \subseteq M$ is said to be \textit{finitely generated}: for some finite subset $A \subset M,\ N = RA$. \\
    A submodule $N \subseteq M$ is called \textit{cyclic} if: there exists $a \in M$ such that $N = R a = \{r \cdot a \mid r \in R\}$.
\end{definition}
